% Frozen preamble for an expanded derivation.
%
% Copied into the output directory alongside the .tex so the artifact still
% builds after this skill is gone -- the same rule _shared/md2pdf states.
%
% Layout note earned from _shared/md2pdf's scars: the per-step block is NOT
% a seven-column table. Wide tables run off the right margin and the text is
% silently clipped. Two bands instead: the move on top, the small print below.
\usepackage[margin=2.4cm]{geometry}
\usepackage{amsmath,amssymb,amsthm,mathtools}
\usepackage{booktabs}
\usepackage{longtable}
\usepackage{array}
\usepackage{xcolor}
\usepackage{tcolorbox}
\usepackage{enumitem}
\usepackage[hidelinks]{hyperref}
\usepackage{microtype}

\tcbuselibrary{breakable,skins}

% The title bar carries the step number and the move name. It was originally
% white-on-light and effectively unreadable in print -- caught by rendering the
% pages and looking at them, which no test would have flagged.
\definecolor{stepframe}{HTML}{7B8794}
\definecolor{steptitle}{HTML}{2E3440}
\definecolor{stepback}{HTML}{FBFCFD}
\definecolor{smallprint}{HTML}{4C566A}
\definecolor{gapblocking}{HTML}{A5303A}
\definecolor{gapsubstantive}{HTML}{B07219}
\definecolor{gapminor}{HTML}{4C566A}

\newtheorem{claimenv}{Claim}
\theoremstyle{definition}

% \stepblock{n}{before}{after}{move}{licensed by}{breaks if}{checked}{gloss}
%
% Band 1 carries the mathematics: what went in, what came out, and the name of
% the move. Band 2 is the small print that makes the step auditable.
\newcommand{\stepblock}[8]{%
  \begin{tcolorbox}[breakable, enhanced, colback=stepback, colframe=stepframe,
                    colbacktitle=stepframe!22!white, coltitle=steptitle,
                    boxrule=0.4pt, left=6pt, right=6pt, top=5pt, bottom=5pt,
                    title={\small\bfseries Step #1\hfill\normalfont\ttfamily\scriptsize #4}]
  {\setlength{\abovedisplayskip}{2pt}\setlength{\belowdisplayskip}{2pt}%
   \begin{flushleft}
   \textcolor{smallprint}{\scriptsize FROM}\\[1pt]
   $\displaystyle #2$\\[3pt]
   \textcolor{smallprint}{\scriptsize TO}\\[1pt]
   $\displaystyle #3$
   \end{flushleft}}
  \vspace{2pt}\hrule\vspace{4pt}
  {\footnotesize\color{smallprint}
   \begin{description}[leftmargin=6.4em, style=sameline, itemsep=1pt, topsep=0pt]
     \item[Licensed by] #5
     \item[Breaks if] #6
     \item[Checked] #7
     \item[In words] #8
   \end{description}}
  \end{tcolorbox}
}

% A step the expansion could not complete. Rendered inline, where the step would
% have been, because a gap is part of the derivation and not an appendix to it.
\newcommand{\stepgap}[3]{%
  \begin{tcolorbox}[breakable, enhanced, colback=white, colframe=gapblocking,
                    boxrule=0.9pt, left=6pt, right=6pt, top=5pt, bottom=5pt,
                    title={\small\bfseries #1\hfill\normalfont\scriptsize could not be made explicit}]
  {\footnotesize\textbf{What is missing.} #2\par\vspace{3pt}
   \textbf{What would close it.} #3}
  \end{tcolorbox}
}

\newcommand{\notationrow}[3]{#1 & #2 & #3 \\}
